% readme.tex -- a short example of how each Stata Journal insert should be
% organized.

\inserttype[st0001]{article}
\author{Michler, Josephson, Barriga-Cabanillas, Michuda, Oliver}{%
  Jeffrey D. Michler\\University of Arizona\\Tucson, AZ\\jdmichler@arizona.edu
  \and
  Anna Josephson\\University of Arizona\\Tucson, AZ\\jdmichler@arizona.edu
  \and
  Oscar Barriga-Cabanillas\\World Bank\\Washington, D.C.\\obarriga@worldbank.org
  \and
  Aleksandr Michuda\\Swarthmore College\\Swarthmore, PA\\amichud1@swarthmore.edu
  \and
  Jeffrey C. Oliver\\University of Arizona\\Tucson, AZ\\jcoliver@arizona.edu
}
\title[wxsum]{wxsum: A command for processing temperature and precipitation data}
\maketitle

\begin{abstract}
Weather data are now a standard input in socioeconomic analysis, yet computing seasonal weather metrics from daily data remains error-prone and time-consuming. In this article, we introduce \texttt{wxsum}, a community-contributed command that reduces this cost. The command takes daily precipitation or temperature data and calculates 18 precipitation and 13 temperature measures, ranging from summary statistics (mean, standard deviation) to shock measures (z-scores, dry spells, killing degree days). By standardizing these calculations, \texttt{wxsum} lowers the probability of coding errors and allows researchers to move quickly from raw weather data to analysis-ready variables.

\keywords{\inserttag, wxsum, Earth observation, weather data, data management, summary statistics}
\end{abstract}

% discussion of the Stata Journal document class.
\input wxsum_body_v1.tex
% discussion of the Stata Press LaTeX package for Stata output.
%\input stata.tex

\bibliographystyle{sj}
\bibliography{wxsum_ref_v1}

\section*{Acknowledgments}
We thank Brian McGreal for help in the early debugging of the \texttt{wxsum} command and Emanuele Clemente for final testing. Generative artificial intelligence tools were used during the final debugging and manuscript preparation process. Specifically, Anthropic's Claude Opus 4.6 Thinking, OpenAI's ChatGPT 5.5 Thinking, and Google Lab's Jules based on Gemini 3.1 Pro were used in final debugging, testing, and validation. The same Claude and GPT models were used to edit and format the article's text. The authors wrote, reviewed, and are responsible for the final content of the article and software.

\begin{aboutauthors}
Jeffrey D. Michler is an Associate Professor in the Department of Agricultural and Applied Economics at the University of Arizona.

Anna Josephson is an Associate Professor in the Department of Agricultural and Applied Economics at the University of Arizona.

Aleksandr Michua is a Visiting Assistant Professor at the Department of Economics at Swarthmore College.

Oscar Barriga-Cabanillas is an Economist at the World Bank.

Jeffrey C. Oliver is a Data Science Specialist in the University of Arizona Libraries.

\end{aboutauthors}

\endinput
